\subsection{熟人强奸与"强奸迷思"：为什么大多数受害者认识加害者}

前文介绍了强奸犯罪的法律界定与受害者的应对路径。本节处理另一个决定现实的认知问题：\textbf{人们对强奸的想象，与强奸实际发生的方式相差很远}——而这种差距，正是受害者不被相信、不被支持的原因。

\textbf{一、被数据纠正的想象}
\begin{itemize}
\item \textbf{大多数受害者认识加害者}：陌生人作案是少数，加害者更常是熟人、朋友、同学、同事、约会对象、前伴侣或家庭成员。这直接导致两个后果：受害者在过程中更难识别"这是强奸"，事后更难被他人相信；
\item \textbf{多数情形不涉及明显暴力}：胁迫的手段常是\textbf{施压、灌酒、孤立、利用信任与权力差、反复纠缠、以关系为要挟}，而不是影视作品里的暴力袭击。没有反抗痕迹，是常见情况而不是例外；
\item \textbf{受害者反应并非"战或逃"}：面对威胁，大脑可能进入"僵直"或"讨好"状态——身体动不了、说不出话、配合以图尽快结束。这些反应由神经系统自动触发，与"愿意"无关，却常被误读为"没有反抗就是同意"。
\end{itemize}

\textbf{二、"强奸迷思"逐条澄清}
"强奸迷思"指对强奸与受害者的一组错误而顽固的信念。它之所以重要，是因为这些信念会进入亲友的对话、甚至进入部分办案人员的判断。
\begin{itemize}
\item \textbf{"她穿成那样/喝那么多，是自己招来的"}——责任在实施侵害的人。衣着与饮酒都不是同意，也不能减轻加害者的责任；
\item \textbf{"她没有明确拒绝，甚至没反抗"}——同意需要\textbf{明确且自愿}的表达。没有"不"，不等于"是"；在失去表达能力的状态下（昏迷、严重醉酒、被药物影响），"同意"在法律上无效；
\item \textbf{"他们是情侣/夫妻，不算"}——婚姻与恋爱关系不是身体的无限授权。婚内强奸同样构成犯罪（参见《婚姻内的同意：结婚证不是"永久同意书"》一节）；
\item \textbf{"男性不会被强奸"}——男性同样是受害者，加害者可以是任何性别，也可以是同性；
\item \textbf{"事后还能正常生活，说明没受害"}——创伤反应的形态与时间线差异极大，"看起来没事"不能作为判断标准。
\end{itemize}

\textbf{三、为什么"熟人"这一层让一切更难}
\begin{itemize}
\item \textbf{识别困难}：加害者的行为常被包裹在"追求""暧昧""他喝多了"的叙述里，受害者自己也常在第一反应中否认真实性质；
\item \textbf{社交代价}：共同的朋友圈、工作关系、共同社群会让追责变得昂贵——"闹大了以后怎么见人"，是很多人沉默的真实原因；
\item \textbf{证据更少}：没有陌生人作案那种监控与目击，过程往往发生在私密空间，事后陈述成为核心证据，而"熟人关系"又常被用来��疑陈述的可信度。
\end{itemize}

\textbf{四、可做的事（对当事人与旁观者）}
\begin{itemize}
\item \textbf{对当事人}：把"我是不是想多了"变成"我不舒服"——以自身的感受为判断起点；尽早进行医学处置与证据保全（72 小时内的 PEP、紧急避孕、检测与取证），是否需要报警可以稍后决定（参见上一节）；
\item \textbf{对亲友}：说的第一句话最重要。\textbf{"我相信你""这不是你的错""你想怎么做我都陪着你"}——三句话比任何具体的建议都有用；不追问细节、不替对方做决定；
\item \textbf{对旁观者与社群}：不在群里讨论受害者的私人信息，不转述细节，不传播任何尚未核实的版本；看到"她是不是自愿的"这类讨论时，明确指出责任归属不在受害者；
\item \textbf{对可能存在的施压者}：任何"私了""别闹大""为了大家的面子"的劝说，都在制造二次伤害。正确的做法是把决定权完整地交还给受害者。
\end{itemize}
